\chapter{雷达 SLAM}

\textbf{作者：} Martin Magnusson、Christoffer Heckman、Henrik Andreasson、Ayoung Kim、Timothy Barfoot、Michael Kaess 和 Paul Newman

本章探讨雷达（RAdio Detection and Ranging）在 SLAM 中的应用。与相机和 LiDAR 相比，雷达在 SLAM 中的应用相对不足；但由于它能够在低能见度、远距离条件下工作，并且能够原生产生速度信息，其受欢迎程度正在上升。我们首先讨论机器人中通常使用的雷达传感器类型、其独特的感知原理，以及雷达应用面临的一些挑战（第 9.1 节）。随后讨论雷达滤波、雷达里程计（第 9.2 节）、地点识别（第 9.3 节）以及 SLAM（第 9.4 节）；图 9.1 展示这些部分如何衔接。最后讨论雷达 SLAM 数据集（第 9.5 节），并展望雷达的未来应用（第 9.6 节）。

\begin{figure}[htbp]
\centering
\includegraphics[width=\linewidth,height=0.78\textheight,keepaspectratio]{6485c2c26e82c18f65c65a46aae600bc8df605abf181ff0ef5eb02e52a4bf1d1.jpg}
\caption{图 9.1 基于雷达的 SLAM 信息流与其他 SLAM 系统遵循相同模式，但各过程的细节会针对雷达特性进行调整。}
\label{fig:chapter-9-6485c2c26e82c18f65c65a46aae600bc8df605abf181ff0ef5eb02e52a4bf1d1}
\end{figure}

\section{雷达简介}

\subsection{传感器类型}

下面介绍机器人中最常见的两类雷达：旋转雷达和片上系统（system-on-a-chip，SoC）雷达（见图 9.2）。二者各有优缺点，主要区别在于如何进行“驱动”：旋转雷达机械旋转单个天线，而 SoC 雷达使用多个天线，通过组合各天线信号推断反射发射信号的物体的角度和距离。

\begin{figure}[htbp]
\centering
\includegraphics[width=\linewidth,height=0.78\textheight,keepaspectratio]{63ca43f7732c9806485524adeb61d94b9f31e3230945fe3cf6877774edbcb98e.jpg}
\caption{图 9.2 机器人中使用的两类主要雷达：旋转雷达（左上）和相控阵 SoC 雷达（右上）。图中给出每类传感器的示例及其产生的主要数据产品。原始数据通常随后通过雷达滤波转换为稀疏点云（下方）。右下为带烟雾隧道中的三个点云：红色为受烟雾影响、测程显著受限的 LiDAR 点云；蓝色为二维旋转雷达点云，清晰显示各个方向的隧道壁，且可见较远距离；绿色为 SoC 雷达三维点云，可见车辆前方的墙壁和地面。左下为同一场景的热成像与视觉/RGB 图像数据，以及从中提取蓝色点云的极坐标雷达图。}
\label{fig:chapter-9-63ca43f7732c9806485524adeb61d94b9f31e3230945fe3cf6877774edbcb98e}
\end{figure}

\subsubsection{旋转雷达}

旋转雷达有时也称扫描雷达或成像雷达。它利用旋转雷达传感器构建周围环境的精确极坐标表示，如图 9.2 所示。其主要数据产品是极坐标雷达图。成像雷达的特点是能够检测距离超过 100 m 的物体。在某些模式下，还可以利用多普勒效应确定这些物体的速度。

不同于 LiDAR 系统，旋转雷达只能提供二维平面数据，无法测量所检测物体的高度。即使反射可能来自天线垂直波束宽度内的不同高度，这一限制仍然存在。此外，由于其工作机制是在每个天线角度发射和接收单个脉冲，这些雷达有时无法提供速度数据，因为速度计算需要多个脉冲。较新的工作利用多个相邻方位角的信号恢复速度。

与目前普遍使用的多垂直光束三维 LiDAR 传感器类似，构造机械旋转三维雷达在实践中不可行，因为与激光二极管相比，天线及其聚焦机构尺寸大得多。

\subsubsection{SoC 雷达}

片上系统（SoC）雷达将处理单元集成到数量极少的芯片中；这些芯片直接安装在贴片天线上（发射器/接收器阵列），或集成到印刷电路板中。由于集成架构和没有运动部件，SoC 雷达相较旋转雷达更轻、更省电。SoC 雷达的精度和分辨率取决于天线阵列设计，以及制造商用于整合多个天线测量的专有处理技术。其主要输出是雷达数据立方体（radar datacube）。

SoC 回波通常按方位角、俯仰角和距离映射到球面坐标。由于径向速度增加了一个维度，这类雷达通常称为 3+1D 或 4D。相反，由于垂直排列天线较少而垂直分辨率有限或缺失的系统，称为二维阵列雷达，产生 2+1D 数据产品。

\subsection{雷达感知原理}

本节讨论与机器人应用相关的基本操作、天线和数据类型，以及应用挑战。下面各小节将分别概述它们，并说明雷达与其他测距传感器的差异。

毫米波（millimeter-wave，mmWave）雷达指电磁波长在 1--10 mm、频率范围为 30--300 GHz 的雷达系统。在这一范围内，多数雷达传感器工作于 76--81 GHz 频段，因为高级驾驶辅助系统（ADAS）等汽车应用已在该频段获得专用频谱。

\subsubsection{雷达散射截面}

毫米波雷达通过发射电磁脉冲实现感知；脉冲传播直到遇到物体，再反射回雷达。雷达散射截面（radar cross section，RCS）受物体材料组成、大小和形状影响，在确定每个物体对电磁脉冲的反射强度方面起着关键作用。从本质上说，RCS 表示一个假想球体的大小，该球体反射的能量与目标物体相同。因此，相较于小物体或行人（RCS 较低），车辆和厚混凝土墙等更大、更坚固的结构具有更高 RCS。

雷达强度指物体回波的强度，是雷达发射功率和物体 RCS 的函数。也就是说，该强度主要由雷达发射功率和所遇目标的 RCS 决定。文献表明，这种强度度量对于提取物体语义细节或辅助导航至关重要，因为较强回波往往对应环境中独特且易于识别的特征 \textsuperscript{[1240]}。雷达信号强度还受到多个因素影响，包括所用天线类型、发射电磁脉冲的特性以及天线检测物体回波的能力。

\subsubsection{FMCW 测距}

雷达至少包含一个发射天线（TX）和一个接收天线（RX）。\textsuperscript{1} 在 FMCW 雷达中，TX 天线发射频率稳步增加的射频脉冲，称为 chirp。该过程通常称为锯齿调制；后文还将讨论另一种基于 chirp 的技术，即三角调制。chirp 在环境物体上反射，RX 天线捕获回波（见图 9.3）。接收到 chirp 后，信号先被放大，再与原始 TX chirp 混合或相减，生成中频（intermediate frequency，IF）信号。由于发射和接收信号的斜率相同，混合过程产生频率恒定为 $f _ { 0 }$ 的正弦 IF 信号。雷达系统的性能，包括测距和速度检测能力，会受到多种 chirp 参数影响，例如带宽（初始频率与终止频率之差）、chirp 斜率以及 chirp 之间的持续时间。

\begin{figure}[htbp]
\centering
\includegraphics[width=\linewidth,height=0.78\textheight,keepaspectratio]{ab8cdc8bfb238ce9142f4dad7c7d43d18c7a047bed5573bff0b95eac5e6f28cf.jpg}
\caption{图 9.3 调频连续波（FMCW）雷达通过发射频率递增或递减的 chirp（波），这些波从目标反射后被接收。通过分析反射信号相对于发射信号的频率和相位偏移，可以确定距离（及速度）。两种常用频率调制策略是锯齿调制和三角调制；后者的优势是能够将多普勒频移与距离频移分离。}
\label{fig:chapter-9-ab8cdc8bfb238ce9142f4dad7c7d43d18c7a047bed5573bff0b95eac5e6f28cf}
\end{figure}

这种频率调制的核心思想，是将时间信息编码到连续波上，从而计算距离。与之相反，调幅方法中，返回信号的精度完全取决于信号带宽。频率调制对信噪比和目标 RCS 等问题更具韧性，而这些因素可能影响单次回波的清晰度。在 FMCW 中，脉冲序列中的波以独特时间间隔交错，信号保持相干。这种相干性使每个 chirp 能够准确关联到其发射--接收天线对及其在序列中的位置，从而按幅度和相位对信号进行相关。这相较更早的不相干脉冲式飞行时间雷达系统有显著进步，后者依赖不相干幅度测量，需要熟练操作员从杂波中筛选并分离重要信号。

在 FMCW 中，距离计算依赖信号发射与接收回波之间的时间间隔。利用光速 $c$ 和初始到达时间 $t _ { 0 }$，可按下式计算物体距离 $d$：

\begin{equation}
d = \frac {c}{2} t _ {0 }.\tag{9.1}
\end{equation}

图 9.3 左上部分的绿色虚线显示了生成 IF 信号的区域。发射信号与接收信号之差就是 IF 信号。IF 信号是正弦波 $\sin(2\pi f _ {0} t + \phi _ {0})$，其频率与常数 $f _ { 0 }$ 成正比；该常数从 $t _ { 0 }$ 延续到 $T$，只取决于目标距离，并有相位偏移 $\phi _ { 0 }$。

频率 $f _ { 0 }$ 由目标距离、发射 chirp 的持续时间 $T$ 以及发射信号带宽 $B = f _ { M } - f _ { m }$ 决定。带宽和发射时间通过 chirp 斜率 $S = B / T$ 关联：

\begin{equation}
f _ {0} = \frac {2 B d}{T c} = \frac {2 S d}{c} \quad \Rightarrow \quad d = \frac {c f _ {0}}{2 S},\tag{9.2}
\end{equation}

此时距离 $d$ 表示为所测中频 $f _ { 0 }$ 的函数。

虽然本节在频域中给出公式，但现实世界中的信号采集主要通过数字采样完成。采样由高频模数转换器（analog-to-digital converter，ADC）以很高频率进行，通常每 100 ns 或更短时间采样一次。随后，采样数据依次进行快速傅里叶变换（fast Fourier transform，FFT），生成频率--幅度图，从中可以辨识信号峰值。距离频率通过对单个 chirp 及其对应回波信号数据进行 FFT 得到；速度频率则通过对多个 chirp 和回波序列执行一系列 FFT 计算。

\subsubsection{用锯齿调制确定距离和速度}

锯齿频率调制（见图 9.3）通常用于 SoC 雷达和一些旋转雷达。IF 信号的相位 $\phi _ { 0 }$ 可表示为信号波长 $\lambda$ 和距离 $d$ 的函数：

\begin{equation}
\phi _ {0} = 2 \pi f _ {m} t _ {0} = \frac {4 \pi d}{\lambda}.\tag{9.3}
\end{equation}

尽管基频 $f _ { 0 }$ 和相位 $\phi _ { 0 }$ 都是距离 $d$ 的函数，但相位仅在足够小的距离值下有效，且存在相位缠绕。因此，它通常不用于测距，而用于测量小的距离变化 $\Delta d$；在速度估计中，相位会对这类变化产生线性响应。

为估计径向速度（见图 9.4），如图 9.3 左上所示，至少使用两个连续上行 chirp。序列中每个 chirp 的时间间隔很小，量级约为 40 微秒，因此两次采样的距离测量以及相应 IF $f _ { i }$ 和 $f _ { i+1 }$ 几乎相同。然而，IF 信号的相位不同。这一相位差 $\Delta \phi$ 对应物体运动；估计速度 $v$ 由相位差确定：

\begin{figure}[htbp]
\centering
\includegraphics[width=\linewidth,height=0.78\textheight,keepaspectratio]{a43c559321a1564a744d1f4cac675b0a1abf0ca83c92f632fbc1a8641dbdc68d.jpg}
\caption{图 9.4 雷达可利用多普勒效应测量设备与其成像场景之间的径向速度。}
\label{fig:chapter-9-a43c559321a1564a744d1f4cac675b0a1abf0ca83c92f632fbc1a8641dbdc68d}
\end{figure}

\begin{equation}
\Delta \phi = \frac {4 \pi \Delta d}{\lambda} = \frac {4 \pi v T}{\lambda},\tag{9.4}
\end{equation}

化简为

\begin{equation}
v = \frac {\lambda \Delta \phi}{4 \pi T}.\tag{9.5}
\end{equation}

由于速度是相位的函数，在 $| \Delta \phi | < \pi$ 时可获得无歧义的最大可检测速度。因此，$v _ { \mathrm {max} } = \lambda / (4T)$ 是信号波长和 chirp 间时间间隔的函数。

\subsubsection{用三角调制确定距离和速度}

也可以使用三角频率调制（见图 9.3）；例如，它有时用于旋转雷达。实际上，从上行 chirp 得到中频时，它包含两个分量：(i) 由信号飞行时间造成的频移（前文已讨论）$f _ {0,t}$；(ii) 若雷达和目标存在相对速度，则由多普勒效应造成的表观频移 $f _ {0,d}$：

\begin{equation}
f _ {0, \mathrm{up}} = f _ {0,t} + f _ {0,d}.\tag{9.6}
\end{equation}

但是，若在上行 chirp 后紧接着使用一个从相同目标反射的下行 chirp，时间频移的符号会翻转，而多普勒频移的符号不会：

\begin{equation}
f _ {0, \mathrm{down}} = - f _ {0,t} + f _ {0,d}.\tag{9.7}
\end{equation}

由这两个方程可求出 $f _ {0,t}$ 和 $f _ {0,d}$：

\begin{equation}
f _ {0,t} = \frac {f _ {0, \mathrm{up}} - f _ {0, \mathrm{down}}}{2}, \quad f _ {0,d} = \frac {f _ {0, \mathrm{up}} + f _ {0, \mathrm{down}}}{2}.\tag{9.8}
\end{equation}

\begin{figure}[htbp]
\centering
\includegraphics[width=\linewidth,height=0.78\textheight,keepaspectratio]{94ee3c464fa3e6a26cf02955d1a3bcc364b0f89d85a8af88d806782a2a60bb64.jpg}
\caption{图 9.5 使用相控阵雷达估计到达角。TX 发射天线发射的信号，由相距 $\ell$ 的两个 RX 接收天线接收时所测得的相位差，对应于目标的角度 $\theta$。}
\label{fig:chapter-9-94ee3c464fa3e6a26cf02955d1a3bcc364b0f89d85a8af88d806782a2a60bb64}
\end{figure}

最后，根据这两个分量可计算距离和速度：

\begin{equation}
d = \frac {c f _ {0,t}}{2S}, \quad v = \frac {c f _ {0,d}}{2ST}.\tag{9.9}
\end{equation}

值得注意的是，先前式（9.2）中的距离计算未校正多普勒效应，而这里的计算已校正。使用三角调制的缺点是增加了延迟，因为计算距离和速度时需要稍早的数据；但不再需要处理信号相位。

\subsubsection{确定 SoC 雷达的角度}

对于旋转雷达，在每一时刻波束集中于单一方位方向，因此确定目标角度十分简单。SoC 雷达的角度估计稍复杂，但可用与上述类似的相位差计算实现。

给定多个间距为 $\ell$ 的接收单元，反射中会产生距离差 $\Delta d$。利用几何关系 $\Delta d = \ell \sin \theta$，可由式（9.3）的变体推导到达角 $\theta$。如图 9.5 所示，接收信号到达第二个接收天线必须多走 $\ell \sin \theta$ 的距离。这对应于两个 RX 天线接收信号之间 $\Delta \phi = (2\pi / \lambda) \ell \sin(\theta)$ 的相位差。给定所测相位差 $\Delta \phi$，可计算到达角 $\theta$：

\begin{equation}
\theta = \arcsin \frac {\lambda \Delta \phi}{2 \pi \ell}.\tag{9.10}
\end{equation}

原则上，一个 TX 天线和两个 RX 天线已足以确定目标角度；但多于两个 RX 天线可提高分辨率，因此可以区分相邻的多个目标。返回信号在每个 RX 天线处都会额外偏移 $\Delta \phi$。对 RX 天线上的信号采样并对此信号序列执行 FFT，可以可靠地估计 $\Delta \phi$。

\begin{figure}[htbp]
\centering
\includegraphics[width=\linewidth,height=0.78\textheight,keepaspectratio]{2536714fdcbb28d955c776ed219adb1c6aded4e3f3f7bb16847fa9fe86f0fdd0.jpg}
\caption{图 9.6 (a) 中的极坐标雷达图转换为 (b) 中的欧氏坐标后，可从 (c) 的放大图观察到几类区别于其他传感器的雷达噪声。最常见的是散斑噪声回波，模糊杂波用绿色圈出。多径反射发生在回波接触天线前先从附近墙壁或地面反弹时，从而产生真实目标的反射；一串重复回波用红色圈出。原始图像取自 Mulran 数据集 \textsuperscript{[567]}。}
\label{fig:chapter-9-2536714fdcbb28d955c776ed219adb1c6aded4e3f3f7bb16847fa9fe86f0fdd0}
\end{figure}

上述示例说明了“SIMO”相控阵雷达系统（single input, multiple output，单输入多输出）。多数用于 SLAM 的雷达传感器是具有多个 TX 和 RX 天线的 MIMO（multiple input, multiple output，多输入多输出）系统。只要 RX 天线能够区分多个 TX 天线的信号，与其将 RX 天线数量加倍，不如增加一个 TX 天线也可以达到相同分辨率。可以采用不同技术确保 TX 信号不相关（正交），例如频分（每个发射器使用不同频段）、码分（每个发射器发送由唯一代码序列调制的信号），或时分多址（time-division multiple access，TDMA），即每个发射器使用不同时间槽。同样原理也可应用于能够同时测量方位角和俯仰角的二维 TX--RX 阵列，产生 3+1D 数据。

\subsection{雷达应用的挑战}

与任何传感器系统一样，雷达技术也带来一系列需要在开发中审慎考虑的挑战。这些挑战包括多种在雷达中尤其突出的噪声类型，例如多径反射、因波束宽度较大造成的有偏且稀疏的距离读数、接收机饱和和散斑噪声。图 9.6 给出了部分噪声现象的示意，它显示了转换为欧氏坐标的极坐标雷达图。第 9.1.4 节将讨论一些相关雷达滤波技术。

\subsubsection{散斑噪声}

雷达测量噪声来自多个来源，包括热噪声、电子缺陷以及目标 RCS 的变化。雷达发射电磁脉冲时，捕获天线视场内所有物体反射回来的能量。脉冲与物体的相互作用会散射雷达波，导致相长和相消干涉。这类相互作用无论信号来源如何，都可能产生假信号，或抵消天线接收到的合法回波。这些因素导致频域信号变化，其中最显著峰值混合了真实目标和虚警。缺乏区分真实回波与虚假回波的机制时，传感器会生成散布点的图案，通常称为散斑噪声。为了准确识别对位姿估计和特征匹配至关重要的路标，必须估计这些反射周围的不确定性，可能还需跨多次扫描进行估计。

\subsubsection{多径}

除散斑噪声外，多径回波也是另一种测量误差，由与单一物体相关的不同探测路径产生。设想传感器前方有一个路标：部分发射射线可能直接到达该路标，而其他射线可能先从地面或墙壁反射，之后才到达天线。雷达会将其解释为路标位于道路下方或墙壁后方，从而感知到所谓“鬼影物体”（ghost objects）或静态离群值。消除这些离群值对于确保点云建图或定位的可靠性至关重要。

\subsubsection{运动引起的畸变}

包括 LiDAR 和雷达在内的扫描传感器天然存在运动畸变。旋转雷达尤其如此，它通过单次旋转构建每幅极坐标图像。因此，若传感器正在运动，同一旋转开始和结束时捕获的单个物体位置不同。对低频运行的传感器或车辆快速运动时，这一差异会很显著。例如，常用成像雷达 Navtech CIR 304 的工作频率为 4 Hz，使得原始帧难以准确对齐。还有更快的旋转雷达，如以 10 Hz 旋转、与许多 LiDAR 相近的新型 Navtech RAS3，以及可达 50 Hz 的 Indurad iSDR。即便如此，当传感器安装在汽车等快速移动平台上时，运动引起的畸变仍可能显著。

\subsection{雷达滤波}

假目标的出现和分布（第 9.1.3 节）会随时间变化，且具有不可预测的参数；因此，简单阈值等静态滤波方法不足以应对，因为它们可能让虚警通过。为此，研究者设计了能够动态估计虚警分布的方法。这类滤波方法通常考虑一个方位方向上的测量，试图估计哪些距离单元包含真实目标，哪些属于虚警。

恒虚警率（constant false alarm rate，CFAR）滤波器 \textsuperscript{[811]} 是一种广泛实现的方法，旨在在动态变化且不均匀的干扰中维持指定虚警概率。该过程先将信号（例如在雷达 ADC 样本经 FFT 后得到的频域表示）划分为称为单元的离散分段，再使用滑动窗口评估这些单元。窗口核心是一组待测单元（cells under test，CUT）。通过比较 CUT 的强度与其前后相邻的训练单元强度进行评估。在一些实现中，训练单元和 CUT 之间还会设置保护单元，以防 CUT 的局部影响影响训练单元幅值。根据 CUT 集的强度是否超过计算阈值，决定接受或拒绝；该阈值由周围训练单元（不包括保护单元）的比较强度导出。CFAR 有多种变体，经典版本为单元平均 CFAR（cell-averaging CFAR，CA-CFAR）：阈值相对于训练单元的平均功率计算，并按根据目标虚警概率选择的阈值乘数缩放。常见替代方案为有序统计 CFAR（ordered-statistic CFAR，OS-CFAR），它使用训练单元的中位数等更鲁棒的统计量，而非均值。

图 9.7 展示了 CA-CFAR 阈值如何适应雷达信号，以及与其他两种滤波方法相比，哪些距离单元被选为目标或噪声。

专门针对雷达里程计设计和测试的 CFAR 变体是 BFAR（bounded false alarm rate）\textsuperscript{[25]}。它只需对 CFAR 检测器计算的输出 $Z$ 应用仿射变换 $T = aZ + b$；其中 $b$ 是可学习参数，用于缩放输出，使其在 CFAR 输出与固定阈值之间混合。

尽管 CFAR 及其变体用于许多雷达应用，若干雷达里程计和 SLAM 流水线采用更简单的滤波策略。CFAR 是为目标检测开发的，在该任务中不遗漏弱检测可能很重要；而为里程计过滤雷达时，关注点是仅保留能够随时间和从不同视点可靠检测到的点。一种常用技术是在每个方位角上，选择高于静态阈值的 $k$ 个最强回波，其中 $k$ 从 1 起。也可使用统计阈值，选择强度高于均值一个标准差的所有点。

一些近期方法改用机器学习技术提高雷达输出的准确性和分辨率，通常使用 LiDAR 数据作为训练真值。Cheng 等人 \textsuperscript{[195]} 使用生成对抗网络（generative adversarial network，GAN）根据距离--多普勒速度矩阵生成点云。Xu 等人 \textsuperscript{[1210]} 训练回归器和分类器：回归器输出改进的高分辨率深度读数，分类器估计数据是否超出测程。这些方法旨在学习一种模型，只保留雷达图中对应于 LiDAR 可以检测到的表面的回波。

\begin{figure}[htbp]
\centering
\includegraphics[width=\linewidth,height=0.78\textheight,keepaspectratio]{4c941393053409c9e2146a4dbd1c8e8efd0dafe7ad7c4b5e5b1f770b4ab19ec1.jpg}
\caption{图 9.7 雷达滤波示例，比较 CFAR 滤波器、恒定功率阈值和 $k$ 最强策略。一个方位方向上的功率/距离曲线以灰色绘制。真实目标的回波在图中呈现为尖峰，但也存在若干模糊峰值。CFAR 产生自适应阈值，以黑线绘制；CFAR 报告的检测以蓝色绘制，在本例中包含多个“虚警”。$k$ 最强滤波器（本例 $k=12$）更保守，仅返回主要目标附近的点。图引自 Adolfsson 等人 \textsuperscript{[6]}（©2021 IEEE）。}
\label{fig:chapter-9-4c941393053409c9e2146a4dbd1c8e8efd0dafe7ad7c4b5e5b1f770b4ab19ec1}
\end{figure}

\section{雷达里程计}

雷达里程计的目标是：给定随时间排序的一组雷达读数，估计传感器的自身运动。雷达里程计方法通常处理某种中间表示，例如最近 $N$ 次雷达读数，或持续裁剪的局部地图。其重点是获得局部尺度上的准确位姿估计。即使里程计方法性能良好，如果不显式考虑回环，误差最终也会无界累积。使用旋转雷达的方法每行驶 100 m 可能累积约 1--2\% 的平移漂移。

许多雷达传感器的一个特征，是能够逐点提供“多普勒”速度估计。第 9.2.1 节将说明如何在无对应关系的情况下利用它估计里程计。除基于多普勒的方法外，还经常通过空间对应关系估计两次相邻雷达扫描之间的相对变换，以确定一次扫描的哪些部分也存在于另一次扫描中。给定一组对应关系，可以计算并优化距离度量。根据扫描器类型，获得对应关系的方式不同：旋转雷达通常产生原始信号，可以像第 9.2.2 节那样直接使用，也可以从原始信号提取包含信息的高层特征（第 9.2.3 节），或提取可用于配准的距离点，类似 LiDAR 里程计（第 9.2.4 节）。

\begin{figure}[htbp]
\centering
\includegraphics[width=\linewidth,height=0.78\textheight,keepaspectratio]{0cabd6dce45ccff6946274bb059e9c0bf50fd880e018b61ab8938a6250bd2efa.jpg}
\caption{图 9.8 Oxford Radar Robotcar 数据集 \textsuperscript{[127]} 上的开环雷达里程计示例。使用 CFEAR 方法 \textsuperscript{[8]} 处理 Navtech 二维旋转雷达数据。真值轨迹以蓝色绘制，里程计估计以橙色绘制；从雷达图提取的点目标以灰色绘制。}
\label{fig:chapter-9-0cabd6dce45ccff6946274bb059e9c0bf50fd880e018b61ab8938a6250bd2efa}
\end{figure}

图 9.8 展示了使用二维扫描雷达进行开环雷达里程计的示例。

\subsection{多普勒里程计}

由多普勒测量获得的径向速度可用于直接估计传感器的线速度。Doer 和 Trommer \textsuperscript{[272]} 将多普勒信息与三点 RANSAC 和最小二乘问题结合，用于估计线速度。

然而，单个雷达数据中的逐点速度测量无法直接观测旋转速度分量，除非能对雷达系统运动学模型作出假设，例如已知雷达相对于旋转中心的安装位置并假设不存在侧滑 \textsuperscript{[551, 354]}。因此，仅依赖单个雷达多普勒信息的雷达里程计系统通常还使用 IMU 数据（更具体地说是陀螺仪）。Huang 等人 \textsuperscript{[490]} 将消费级 IMU 与级联 SoC 雷达结合，在多样三维室内空间中实现低漂移。Kubelka 等人 \textsuperscript{[609]} 比较了多种基于配准的 3+1D 雷达方法与无配准的多普勒 + IMU 里程计 \textsuperscript{[272]}，发现无配准方法产生的误差最低，尤其在特征稀疏环境中；据报道，其在 4.5 km 轨迹上的漂移低至 0.3\%。

\begin{figure}[htbp]
\centering
\includegraphics[width=\linewidth,height=0.78\textheight,keepaspectratio]{f11a8bf09a459fdc3f1f348c3533a5411923c85daba70f97e5d696f569278554.jpg}
\caption{图 9.9 雷达惯性速度估计系统的因子图表示。利用多普勒目标和多组 IMU 测量作为约束，联合估计过去 $K$ 个时间步的状态。图改编自 Kramer 等人 \textsuperscript{[606]}。}
\label{fig:chapter-9-f11a8bf09a459fdc3f1f348c3533a5411923c85daba70f97e5d696f569278554}
\end{figure}

Kellner 等人 \textsuperscript{[551]} 展示了：当雷达安装在具有 Ackermann 转向的车辆上，且已知雷达传感器相对于车辆旋转中心的安装位置时，可以从 2+1D 雷达数据估计线速度和旋转速度。Galeote-Luque 等人 \textsuperscript{[354]} 将其扩展到 3+1D 情形，估计五个自由度（三维线运动以及偏航和俯仰旋转，但不包括滚转）。

由于基于多普勒的里程计模态相当专属于雷达方法，下面以 Kramer 等人 \textsuperscript{[606]} 的工作为例说明。

\subsection{示例：多普勒里程计因子表述}

从雷达整合多普勒因子的一种直接方式，是在包含过去 $K$ 次雷达测量的滑动窗口中估计传感器平台的机体坐标系速度。这些速度通过 IMU 的加速度计测量积分相互连接，从而形成准确速度估计的完整系统。系统结构可以用因子图表示：利用多普勒目标和多组 IMU 测量作为约束，联合估计过去 $N$ 个时间步的状态。

IMU 测量（加速度计和陀螺仪）通常同时受到偏置 $\boldsymbol {b}$ 和重力 $\boldsymbol {g}^{w}$ 的影响。由雷达数据得到的速度估计不含偏置，因此可将加速度计偏置纳入状态向量进行补偿。但要补偿重力影响，需要估计 IMU 姿态，具体是由方向四元数 $\boldsymbol {q}_{s}^{w}$ 表示的俯仰和滚转。于是，完整状态向量可表示为

\begin{equation*}
\boldsymbol {x} = \left[ \boldsymbol {v}^{s\top} \quad \boldsymbol {q}_{s}^{w\top} \quad \boldsymbol {b}^{\top} \right] ^ {\top}.
\end{equation*}

雷达惯性自身速度估计被表述为一个优化问题，其代价函数整合约束和测量，以准确估计传感器平台速度：

\begin{equation}
J (\boldsymbol {x}) = \underbrace {\sum_ {k = 1} ^ {K} \sum_ {d \in \mathcal {D} _ {k}} w _ {d,k} e _ {d,k}} _ {\text {多普勒项}} + \underbrace {\sum_ {k = 1} ^ {K - 1} \| \boldsymbol {e} _ {k} \| _ {\boldsymbol {W} _ {k}} ^ {2}} _ {\text {惯性项}}.\tag{9.11}
\end{equation}

其中，$K$ 是估计状态的过去雷达测量数，$\mathcal {D}_{k}$ 是时刻 $k$ 雷达测量返回的目标集合，$e_{d,k}$ 是多普勒速度误差，$\boldsymbol {e}_{k}$ 是 IMU 误差。对 IMU 误差，误差项由协方差矩阵 $\boldsymbol {W}_{k}$ 加权；对多普勒速度测量，则由相应雷达目标的归一化强度加权：

\begin{equation}
w _ {d,k} = \frac {i _ {d,k}}{\sum _ {j \in \mathcal {D}_{k}} i _ {j,k}},\tag{9.12}
\end{equation}

其中 $w_{d,k}$ 是扫描 $\mathcal {D}_{k}$ 中目标 $d$ 的权重，$i_{d,k}$ 是目标 $d$ 的强度。

本例中的雷达测量由目标集合组成。对每个 $d \in \mathcal {D}$，测量 $[r\quad v\quad \theta\quad \phi]^{\top}$，分别表示目标 $d$ 的距离、多普勒（径向）速度、方位角和俯仰角。多普勒速度测量 $v$ 等于目标与传感器之间相对速度向量 $\boldsymbol {v}^{s}$ 在传感器坐标系中，沿传感器原点到目标的射线 $\boldsymbol {r}^{s}$ 的投影模长，即目标在传感器坐标系中的速度与从传感器指向目标的单位向量的点积：

\begin{equation}
\boldsymbol {v}^{s\top} \left(\frac {\boldsymbol {r}^{s}}{\|\boldsymbol {r}^{s}\|}\right).\tag{9.13}
\end{equation}

该方法假设场景目标静止，仅传感器平台运动。此时，每个雷达目标都可为机体坐标系中传感器组速度估计提供约束。每个雷达目标的速度误差为

\begin{equation}
e _ {d,k} = v _ {d,k} - \boldsymbol {v} _ {k}^{s\top} \left(\frac {\boldsymbol {r}_{d,k}^{s}}{\|\boldsymbol {r}_{d,k}^{s}\|}\right),\tag{9.14}
\end{equation}

其中 $v_{d,k}$ 是时刻 $k$ 目标 $d$ 的径向速度测量，$\boldsymbol {v}_{k}^{s}$ 是时刻 $k$ 的传感器速度，$\boldsymbol {r}_{d,k}^{s}$ 是时刻 $k$ 从传感器指向目标 $d$ 的向量。如前所述，雷达测量受到非高斯噪声影响，雷达扫描还常包含虚假目标数据。可以对多普勒残差使用 Cauchy 鲁棒范数应对这些挑战。此外，本例假定雷达传感器与 IMU 位置和姿态相同；一般并非如此，\textsuperscript{[605]} 给出了一般推导，阐明机体角速度与雷达传感器坐标系速度之间的耦合。

\subsection{直接里程计}

直接处理原始雷达图（如图 9.6 和图 9.2 所示），而不是处理从雷达图中过滤得到的点的方法，称为直接法。

直接方法使用相位相关和 Fourier--Mellin 变换 \textsuperscript{[176, 848]} 等经典信号处理技术。给定两个连续极坐标雷达图，可在极坐标系中通过平移找到相对旋转，其中垂直平移对应方位角变化。利用相位相关，可从使两幅极坐标图一致性最大的像素平移选择相对方向。随后，在笛卡尔坐标系中类似地计算图像相关性，细化平移（图 9.6）。这些直接相关方法假设特定位置的功率回波随时间保持静止，从而能够进行有意义的相关。但这一假设常常失败，尤其是在存在动态物体或噪声伪影常见的雷达数据中。

Barnes 等人 \textsuperscript{[53]} 的“通过运动进行掩膜”（Masking by Moving）方法也使用相关性。它在规则的位姿候选网格上，计算当前扫描与上一扫描旋转副本之间的二维相关。Barnes 等人通过训练卷积神经网络（CNN）避免纳入由噪声造成的虚假特征，以处理功率回波非静止问题。CNN 训练目标是预测一个掩膜，仅保留静止、因而更适合相关扫描匹配的雷达图部分。

与使用全部或几乎全部雷达图的直接里程计方法相反，本章其余方法都是间接方法：它们先选择特定特征或关键点，再在更稀疏的点上运行，以估计里程计。

\subsection{基于特征的里程计}

二维旋转 FMCW 雷达的雷达图本质上是鸟瞰图像，因此自然有许多工作使用计算机视觉领域的基于图像的特征提取和匹配技术（例如 SIFT \textsuperscript{[140, 650]}、SURF \textsuperscript{[473]} 和 ORB \textsuperscript{[546]} 特征）寻找对应关系并估计里程计。Callmer 等人 \textsuperscript{[140]} 匹配群岛中的大尺度岛屿特征，Li 等人 \textsuperscript{[650]} 从卫星雷达中提取特征。与相机图像相比，雷达数据噪声更高且高度依赖环境，因此提取可用于特征匹配的描述子更困难 \textsuperscript{[474]}。与其他传感器相比，雷达特征描述子还可能具有更强地点依赖性，使传感器位姿不同时数据关联更困难（另见第 9.3 节）。FSCD 和 BASD \textsuperscript{[908, 983]} 是专为雷达图设计的关键点提取器和特征描述子示例。

上述技术提取一组显著关键点（可在后续帧中可靠检测），并计算描述关键点周围区域的特征描述子。给定提取的特征集合，通常通过特征描述子匹配计算对应关系，并结合 RANSAC 等鲁棒估计器。Lim 等人 \textsuperscript{[663]} 给出了基于鲁棒估计器的抗离群基于特征雷达里程计示例。给定对应关系后，最小化特征之间的空间距离，通常使用奇异值分解（singular value decomposition，SVD）求最小二乘解。一般而言，相比下一节讨论的基于配准的方法，基于特征的方法对较大初始误差更稳定，因为特征描述子能够鲁棒关联，而配准方法主要依赖点邻近性进行数据关联。

除上述手工设计的特征描述子外，也可通过深度神经网络生成关键点和特征描述子，从而自动生成特征 \textsuperscript{[52]}。其缺点是需要包含雷达图和相似环境真值位姿的训练数据。

\subsection{基于配准的里程计}

LiDAR 测距传感器的点云里程计估计已有成熟方法，相似技术也可用于从雷达图或雷达数据立方体提取的点云。本节说明如何计算这类雷达点云并用于里程计估计。

对图 9.6 所示能够提供原始信号数据的旋转雷达，首先需要通过滤除不对应相关峰值的信号回波，选择要使用的点。\textsuperscript{2} 如第 9.1.4 节所述，许多方法在每个方位角上提取一组距离读数，例如使用 CFAR（图 9.7、第 9.1.4 节）、使用噪声统计量去除冗余或噪声读数 \textsuperscript{[158]}、BFAR \textsuperscript{[25]}，或简单选择每个方位角上最强的 $k$ 个回波。Kellner 等人 \textsuperscript{[550]} 是例外，他们使用 DBSCAN 聚类，同时考虑相邻方位角，而非每次限制在一个方位维度上搜索目标。近期方法还常使用机器学习训练模型，从雷达图估计类似 LiDAR 的点云 \textsuperscript{[195, 1210]}。关键挑战是提取适当数量的读数：太少会丢失相关信息，太多会包含噪声 \textsuperscript{[158]}。例如，研究发现 CFAR 在这方面难以调节 \textsuperscript{[127]}。

使用上述方法之一提取点云后，可以直接使用，也可以通过分析局部邻域估计法向量、平面和点分布等附加信息。雷达数据的配准方法通常类似 LiDAR 扫描配准（见第 8.2 节）；但雷达数据的噪声水平和稀疏性使两两配准更具挑战。

基于配准的雷达里程计常见策略是将新扫描配准到多次过去扫描，过去扫描可以聚合为子地图，也可以作为独立点云集合。与多次扫描配准能够补偿前述雷达数据的若干挑战。加入更多关键帧后，里程计估计对稀疏且含噪的雷达点云不那么敏感。更多对应关系增加了约束，能够减少特征贫乏环境中的漂移；同时也提供时间冗余，即使用多个关键帧时，突然遮挡或移动物体产生的伪对应关系对里程计估计影响较小。

基于配准的雷达里程计方法示例包括连续时间 ICP \textsuperscript{[127]}、功率平移 NDT \textsuperscript{[615]} 和 CFEAR \textsuperscript{[8]}（不要与第 9.1.4 节讨论的 CFAR 滤波方法混淆），它们都使用子地图或多个关键帧。在 CFEAR 中，通过选择每个方位角上最强的 $k$ 个回波，从雷达图提取点云。每个新点云使用点到点、点到线或点到分布误差度量，联合配准到最近的 $s$ 个关键帧，这类似 NDT 扫描配准（见第一部分）。对每个点，在一定半径内邻近点的协方差矩阵估计法向量。根据法向量一致性、平面性（协方差矩阵的条件数）和邻域点数对点对应关系加权。Kung 等人 \textsuperscript{[615]} 使用固定阈值从雷达图提取点云，并利用 NDT 表示将多个点云聚合到雷达子地图中，其中每个点的贡献由其回波信号强度加权。Burnett 等人 \textsuperscript{[127]} 使用 BFAR \textsuperscript{[25]} 从雷达图提取点云并聚合到局部子地图，随后使用连续时间 ICP 表述优化轨迹估计，其中点关联到高斯过程运动先验。

上述方法均基于二维雷达数据。近期基于配准的 3+1D 雷达里程计流水线往往采用相似策略，但点云提取在传感器上完成，因此选择哪些信号峰值作为有效目标，取决于传感器特定阈值，而非显式特征提取。由于 3+1D 点云提供逐点多普勒速度信息，这些方法通常先根据多普勒数据采用最小二乘估计自身速度（见第 9.2.1 节），再进行点云配准予以细化。4DRadarSLAM \textsuperscript{[1271]} 中的基于配准里程计组件使用适配雷达点云的 GICP \textsuperscript{[985]} 变体，以协方差矩阵对点加权；由于方位角和俯仰角精度有限，距离传感器越远的点不确定性越高。4D iRIOM \textsuperscript{[1310]} 采用一对多的分布到分布匹配，以缓解雷达点云噪声和稀疏性：不将每个点匹配到局部子地图中最近的对应分布，而是匹配到最近分布的加权集合。\textsuperscript{[1310, 1271]} 中给出的完整 SLAM 流水线将在第 9.4.2 节进一步介绍。EFEAR-4D 方法 \textsuperscript{[1199]} 将 CFEAR 二维里程计 \textsuperscript{[8]} 扩展到 3+1D 雷达点云。它先计算基于多普勒的自身速度估计，去除与该最小二乘估计不一致、因而可假定来自移动障碍物的离群点；其余配准方案与 CFEAR 类似，即将扫描配准到一系列先前关键帧，并利用法向量一致性和平面性关联和加权单个点匹配。

\subsection{运动补偿}

如第 9.1.3.3 节所述，里程计估计中必须补偿运动畸变。对频率为 4 Hz 的低速旋转雷达，据报道使用恒速度模型进行自身运动补偿可将绝对轨迹误差（absolute trajectory error，ATE）降低 29\% \textsuperscript{[8]}。给定两次连续雷达扫描的时间戳及使用上述方法计算的相对位姿，可以计算速度；由于逐点时间信息也可获得，便可相应移动每个雷达点。以这种方式偏移单次雷达扫描中的点，能够补偿扫描雷达缓慢扫掠期间产生的显著运动。同一模型还用于为刚性扫描配准提供初始估计；每次将运动补偿后的扫描配准到上一扫描后，将其作为节点加入 SLAM 位姿图。

不过，该方法仍在离散位姿图设置中运行。连续时间轨迹表示有利于获得平滑且准确的轨迹，因为可以查询每次传感器读数时刻用于去畸变的位姿。Ng 等人 \textsuperscript{[803]} 在使用汽车 SoC 雷达的流水线中实现了样条表示。Burnett 等人 \textsuperscript{[129]} 使用高斯过程构造因子图，将 IMU 传感器与旋转雷达结合表示轨迹。

\begin{figure}[htbp]
\centering
\includegraphics[width=\linewidth,height=0.78\textheight,keepaspectratio]{9b9170bb16405d02fa667a5307686fd19a220b65332e68c75a974685f15e8871.jpg}
\caption{图 9.10 在同一地点获得的两次雷达扫描。噪声模式造成视觉混叠，使地点识别过程复杂化。}
\label{fig:chapter-9-9b9170bb16405d02fa667a5307686fd19a220b65332e68c75a974685f15e8871}
\end{figure}

\section{雷达地点识别}

与其他传感器模态一样，地点识别（place recognition，PR）是雷达 SLAM 的必要模块。关于一般问题的良好概述可参见第 8 章。如视觉和 LiDAR 地点识别章节所述，在平移和旋转变化下保持不变性至关重要，例如在不同车道行驶，或从另一条街道到达同一交叉路口时。

\subsection{雷达地点识别的独特挑战}

原则上，前面章节讨论的若干地点识别方法，在适当转换数据格式后也可用于雷达数据；但雷达仍具有一些独特特征和挑战（另见第 9.1.3 节）。

雷达地点识别的主要挑战由多个因素造成。首先，雷达数据分辨率低，细节较少，因而可用于识别的可区分特征更少。雷达数据的宽波束导致角度歧义，使远处物体呈现为宽斑块，而从近处观察时其形状会明显不同。此外，较低的信噪比也会给地点识别带来挑战，除非进行充分滤波，如图 9.6 所示。接收机饱和可能产生强径向特征，该特征随目标观测角度显著变化，如图 9.10 样例所示。因此，由于这些挑战，地点的外观可能从附近不同位置观察时发生显著变化。

由于雷达类型不同，地点识别也需要根据传感器类型区别处理。长距离 360° 旋转雷达与测程和视场有限的 SoC 雷达，对地点的感知和描述有显著差异。例如，旋转雷达产生二维雷达图，可将其视为图像，自然可以将受视觉地点识别启发的方法应用于旋转雷达。对雷达图应用 CFAR 等目标检测算法会得到点云，之后可以改用二维 LiDAR 地点识别方法。旋转雷达往往测程很长，对地点识别十分有利 \textsuperscript{[567]}；但其扫描速率较慢，即使车辆低速行驶，也可能产生运动畸变。另一方面，SoC 雷达无需机械旋转天线即可覆盖视场，因此更新速率更快，较少受运动畸变影响。多数使用 SoC 雷达数据进行地点识别的方法处理二维/三维点云数据，而非完整数据立方体。因此，SoC 雷达地点识别通常建立在现有 LiDAR 地点识别方法之上；但 SoC 雷达测量通常受限于较小视场，并且稀疏性和噪声更高。

目前，使用旋转雷达的地点识别文献多于使用 SoC 雷达的文献。部分原因是旋转雷达在更宽视场内捕获更丰富信息，能够更好地捕获周围结构，从而实现鲁棒地点识别，尤其适合室外应用；另一原因是大规模数据集更易获得。但这一优势并不妨碍 SoC 雷达用于地点识别。旋转雷达提供更大视场，但其投影模型仅提供二维信息，丢失高度细节。对于室内地点识别，较短测程和较小视场通常已足够，此时 SoC 雷达的 3+1D 测量仍具有优势。

\subsection{基于学习的雷达地点识别}

对于旋转雷达，可将二维雷达图视为图像，利用二维图像检索进行地点识别。\textsuperscript{3}

Saftescu 等人 \textsuperscript{[963]} 的工作是二维 FMCW 雷达地点识别的早期方法之一。它使用 CNN 表示雷达图，并针对其极坐标性质使用柱面卷积，学习旋转不变表示。随后，可将查询雷达图与参考图像数据库匹配，得到仅基于外观的拓扑地点识别系统。De Martini 等人 \textsuperscript{[246]} 加入位姿细化，得到拓扑度量建图和定位系统。

地点识别方法可使用输入数据的增强实例训练，以提高对小变化的鲁棒性。地点识别数据通常在驾驶沿途按扫描连续记录，因此可以通过采样时间上相邻的帧获得“增强”实例，并通过平移极坐标雷达图人为加入旋转。与这些增强实例形成对比，网络训练不仅要识别相似实例，还要区分时间上相距较远的实例（及其增强版本）。这样，无需知道雷达数据的真实度量位置，也可通过无监督方式获得训练地点识别算法的数据 \textsuperscript{[352]}。

上述方法专为提供 360° 长距离覆盖的旋转雷达设计。SoC 雷达没有运动部件、体积更小，对汽车应用更有吸引力；但其测程较短、视场较小，因此需要略有不同的处理。

Cai 等人 \textsuperscript{[139]} 在将点云投影到二维图像平面后，使用深度时空编码器生成特征向量，作为拓扑地点识别的地点描述子。根据 RCS 对这些向量重新排序后，将其传入 NetVLAD \textsuperscript{[35]} 层，以滤除无关静态特征。在该工作中，车辆周围安装多个雷达传感器，以克服视场有限问题。

Herraez 等人 \textsuperscript{[454]} 展示了单次扫描雷达地点识别流水线：通过学习关注对地点识别重要的显著点，处理数据稀疏和噪声问题。与 Cai 等人 \textsuperscript{[139]} 类似，Herraez 等人也利用 NetVLAD 生成特征编码。但他们使用刚性核点卷积捕获三维上下文信息，而非将三维点云投影到二维图像。进一步地，“点重要性估计器”输出点对地点识别重要的概率。该估计器使用已知对应点云集合训练，在另一扫描中小半径内存在对应关系的查询点被标为重要。RCS 信息通过独立网络纳入，该网络将点的 RCS 数据编码为更紧凑的特征表示。Peng 等人 \textsuperscript{[858]} 以类似第 9.2.1 节多普勒里程计的方法过滤噪声点：使用 RANSAC 估计自身速度，得到可能静止的内点集合，再将剩余离群点作为噪声滤除。Peng 等人没有使用为视觉地点识别设计的 NetVLAD，而是展示了名为 MinkLoc4D 的雷达专用特征提取骨干网络，该网络借鉴 LiDAR 地点识别架构 \textsuperscript{[1320]}。

\subsection{基于描述子的雷达地点识别}

使用手工设计的描述子而非学习得到的嵌入通常同样有效。雷达地点识别的一种常见做法，是直接采用 LiDAR 描述子，例如 Scan Context \textsuperscript{[568]}、RING \textsuperscript{[1212]} 或 M2DP \textsuperscript{[445]}；不过，由于传感器数据格式不同，需要进行适当调整。例如，旋转雷达数据缺少三维结构信息（如高度），但可以用 RCS 代替。此外，还应针对雷达传感器数据的固有挑战采取额外措施，通常需要仔细进行噪声滤波、稀疏性处理和运动补偿。

Hong 等人 \textsuperscript{[474]} 通过将原本为三维点云设计的 M2DP 描述子 \textsuperscript{[445]} 适配到从雷达图提取的二维点云，实现地点识别。此外，他们研究二维平面上的点分布，以判断点云是否可能具有区分性。对二维点执行 PCA 可得到两个特征值，它们描述点沿各特征向量方向的展开程度。当两个特征值差异显著时，说明扫描在某一方向上缺乏特征（例如高速公路行驶数据）；Hong 等人不会将此类扫描用于地点识别。

Jang 等人 \textsuperscript{[510]} 修改了 RING 描述子 \textsuperscript{[1212]} 以适用于雷达。RING 生成的描述子为正弦图形式，具有旋转和平移不变性。两个正弦图之间的相关性应能给出地点识别问题的正确匹配；但雷达图像中的高水平噪声可能使直接比较失效。研究表明，进一步加入自相关能够提高雷达图像的地点识别性能。

Adolfsson 等人 \textsuperscript{[9]} 以多种方式调整由二维雷达图构成的 Scan Context 描述子。首先，他们计算 Scan Context 网格单元中所有点的强度和，以同时编码强度和点密度。其次，每个描述子都由一组经过噪声滤波和运动补偿后聚合的极坐标图像生成，从而减轻第 9.3.1 节所列的一些挑战。仅保留每个方位方向上最强的 $k$ 个回波是一种保守滤波器，往往可以去除雷达点云中原本存在的大部分噪声。在多个已配准点云上生成 Scan Context 描述子，可以进一步弥补这种保守滤波留下的稀疏数据。在行驶过程中使用旋转雷达时，通过恒加速度模型对点云去斜，对于生成可比较的描述子十分重要。

面向 3+1D SoC 雷达的地点识别方法也通常实现 Scan Context 描述子的变体 \textsuperscript{[1271, 1310, 658]}。然而，虚假雷达点很容易扭曲高度测量，因此，对于雷达点云而言，使用每个径向网格单元的最大高度（原始 Scan Context 的做法）并不总是像 LiDAR 那样有效。Intensity Scan Context 描述子 \textsuperscript{[1151]} 是一种替代方案：它在 Scan Context 网格单元中存储点的最大测量强度，而不是高度。另一种方法是在网格单元内使用强度和，从而同时利用强度和点密度；这样既能在描述子中编码垂直结构，又不易受到噪声点位置的影响。使用 3+1D 雷达时，另一个重要因素是：与 Scan Context 所针对的 360° LiDAR 相比，修改后的雷达描述子必须适应窄得多的视场。由于同一地点从两个不同视点观察时可能呈现完全不同的外观，当传感器只覆盖小视场时，很难实现旋转不变性。一种解决办法是根据当前里程计估计进行回环预过滤，只尝试将当前描述子与偏航角处于某一范围（例如 20°）内的帧匹配 \textsuperscript{[1271]}。

\section{雷达 SLAM}

雷达 SLAM 系统一般采用与 LiDAR 或相机 SLAM 解决方案相同的总体结构。本节简要介绍过去二十年中的若干代表性系统，重点关注针对雷达数据进行 SLAM 定制时的特有方面。第 9.4.3 节还将介绍与其他外感知传感器组合的多模态系统。

\subsection{地图表示}

许多雷达 SLAM 系统使用第 5 章讨论过的类似地图表示。然而，雷达测量稀疏且常包含虚假数据，给建图过程带来了独特挑战。一些较早的雷达建图工作直接基于现有目标检测模型建立地图表示，通常使用路标地图，其中每个检测到的目标构成地图的一部分 \textsuperscript{[224, 271, 140, 983]}。这种目标检测也可用于占据栅格地图，例如 Mullane 等人 \textsuperscript{[781]} 的系列工作；他们在建图阶段利用了检测概率 \textsuperscript{[782]}。

由信号中的单个峰表示的检测目标，也可以像 LiDAR 扫描中的点一样处理。例如，这些点可用于创建二维占据栅格地图 \textsuperscript{[730, 762]} 或点云地图 \textsuperscript{[474]}。提取的点云还可以加入额外信息，例如 \textsuperscript{[9]} 计算周围点的分布以估计方向和权重，其方式类似 surfel 或 NDT 单元。直接点云表示在新兴的高分辨率 SoC 雷达中也很常见；这类雷达具有更大的垂直视场（3+1D）和更多 TX/RX 天线（例如 Sensrad Hugin 雷达具有 48 个 TX 和 48 个 RX 天线）\textsuperscript{[1271, 1310]}。

虽然不太常见，但完整的雷达热图（峰值检测前、每个方向和距离网格单元的强度样本）也可用于提供稠密栅格地图 \textsuperscript{[950]}。此时，地图表示环境中每个位置的反射功率全局热图，而非占据证据。随后可通过相关性实现这些热图之间的二维对齐。

Kramer 和 Heckman \textsuperscript{[604]} 在生成里程计的同时，提出了一种用于雷达数据体素建图的新型传感器模型，即使存在视觉遮挡也能创建稀疏地图。该传感器模型采用 Octomap \textsuperscript{[478]} 中用于估计占据单元与自由单元的对数几率方法，但替换了其中的射线投射模型。Octomap 射线投射模型假定传感器首次接触的位置是唯一相关的占据点，而其推广模型利用雷达能够穿透某些材料的能力，在传感器视场内更新体素概率：有雷达回波的单元概率增加，未检测到的扫描单元概率降低，同时不假设射线上的信息。

Nuss 等人 \textsuperscript{[815]} 提出了一种专门考虑雷达的相关栅格地图表示。他们设计了用于处理栅格地图中动态障碍物的状态估计滤波器，称为概率假设密度多实例 Bernoulli 滤波器。该滤波器将栅格单元视为有限随机集合，并动态融合雷达和 LiDAR 数据。

最后，可以使用机器学习技术处理雷达数据密度低、噪声高的问题 \textsuperscript{[1210, 773]}，以 LiDAR 数据作为真值，提高雷达输出的分辨率并降低噪声。Mopidevi 等人 \textsuperscript{[773]} 构建全局雷达地图，并使用预测网络对地图块进行上采样：网络过滤自由空间区域中的噪声，填充稀疏和空缺区域，生成更接近 LiDAR 结果的地图。

近年来，最初为 RGB 数据开发 \textsuperscript{[765]}、后来应用于 LiDAR 数据 \textsuperscript{[1293, 1056]} 的神经场也被用于二维雷达数据 \textsuperscript{[98]}。神经表示的关键特征是对环境进行隐式表示，使神经网络能够接收空间中的一点并返回距离最近表面的距离、颜色和不透明度等量。在 Borts 等人 \textsuperscript{[98]} 的雷达场中，使用物理先验的雷达传感器模型创建隐式神经几何和反射率模型，随后可利用它从未观测视点合成雷达测量。雷达传感器处接收的功率取决于已知的发射功率和天线增益，也取决于由物体大小、雷达反射率和方向性共同决定的 RCS。神经表示学习将测得的 RCS 分解为大小（面积）与反射率和方向性的乘积两部分。

\subsection{雷达 SLAM 流水线}

前几节介绍了基于图的雷达 SLAM 框架的主要组成部分：开环里程计估计、用于检测回环的地点识别以及地图表示，此外还介绍了与雷达 SLAM 相关的物理和技术原理。

本节回顾若干仅使用雷达作为外感知传感器的完整 SLAM 流水线，并讨论它们如何实现这些组成部分。这些流水线通常遵循图 9.1 所示的架构：前端先将原始雷达数据滤波为点云，并通过地点识别模块确定回环、向底层位姿图添加相应约束；后端地图模块执行逐帧里程计，SLAM 模块则在检测到回环后对地图进行全局优化。

\begin{figure}[htbp]
\centering
\includegraphics[width=\linewidth,height=0.78\textheight,keepaspectratio]{9a32bcfd51412b64c55ffd3217286ac6c3d0bd86c4a1ea92a83e0c1c0d314fa4.jpg}
\caption{图 9.11 使用 3+1D 雷达（Sensrad Hugin）和 IMU 输入通过雷达 SLAM 生成的三维地图示例。颜色表示高度。左：回环闭合前；注意左侧三维地图中明显累积的水平和垂直漂移。右：回环闭合后。}
\label{fig:chapter-9-9a32bcfd51412b64c55ffd3217286ac6c3d0bd86c4a1ea92a83e0c1c0d314fa4}
\end{figure}

对于旋转雷达产生的二维雷达图，TBV-SLAM（“trust but verify”，信任但验证）流水线 \textsuperscript{[9]} 建立在第 9.2 节介绍的 CFEAR 二维雷达里程计方法之上。系统按顺序使用每次雷达扫描跟踪传感器位姿，但出于效率考虑，只将较稀疏的关键帧集合加入 SLAM 位姿图。当估计行驶距离超过某个阈值（例如 1.5 m）时，向位姿图添加新的关键帧，并根据其与最新关键帧的对齐创建里程计约束。与通常情况一样，图中的约束需要两个节点之间的相对位姿偏移，以及用协方差矩阵表示的相关不确定性。有研究表明 \textsuperscript{[9]}，使用预先定义且元素较小的对角协方差矩阵，效果优于根据配准代价函数 Hessian 估计协方差。后者原本应能更好地捕获输入点云形状导致的不确定性，例如隧道中的一对点云应在隧道方向上具有更大的不确定性。然而，使用 Hessian 往往会低估或高估不确定性，可能使后端优化轻微错置关键帧。TBV-SLAM 流水线的核心是地点识别模块：它检索多个候选回环（“信任”），随后逐一测试并选出经过验证的最佳候选。Scan Context 描述子 \textsuperscript{[563]}（见第 9.3.3 节）经过调整，同时考虑点密度和信号强度，以替代二维雷达中不可用的高度数据。此外，系统还实现了若干检索和验证回环候选的技术。对于每个关键帧，系统通过对点云施加横向平移偏移，创建多个增强描述子。搜索回环时，将查询描述子与数据库中的所有增强描述子匹配，以处理车辆在不同车道行驶时形成的回环。通过这种方式找到的候选回环随后依据里程计估计进行滤波，而这依赖于长距离内足够准确的里程计。传播查询帧和候选帧之间里程计约束的不确定性后，可以丢弃估计位置相距较远的候选回环。不过，同时考虑描述子相似度和里程计不确定性可以进一步提高回环检索的鲁棒性。也就是说，不是先寻找与查询描述子 $q$ 最相似的候选 $c$，满足 $c = \operatorname*{arg\,min}_{c} d_{\mathrm{descriptor}}(q,c)$，再依据里程计滤波，而是通过

\begin{equation*}
c = \operatorname*{arg\,min}_{c}\left[d_{\mathrm{descriptor}}(q,c)+d_{\mathrm{odometry}}(q,c)\right],
\end{equation*}

直接选择候选。其中，$d_{\mathrm{odometry}}$ 计算的是：在考虑累积里程计不确定性的情况下，帧 $q$ 与 $c$ 位于同一地点的可能性。按此方式找到的雷达扫描对先由 CFEAR 配准模块对齐，最后通过对齐验证模块进行验证；该模块包含重叠度量以及用于检测轻微错位的 CorAl \textsuperscript{[7]} 度量。总体而言，该流水线展示了如何将基于描述子的地点识别适配到雷达数据（考虑二维雷达图中的多重信号回波，并补偿相对稀疏且缓慢的扫描），以及如何筛选多个回环候选，根据几何对齐和前端位姿估计验证最佳候选。

Hong 等人 \textsuperscript{[474]} 的 RadarSLAM 流水线是二维旋转雷达 SLAM 的另一个代表性例子。该流水线直接跟踪雷达图中检测到的关键点来实现里程计（开环位姿跟踪）。系统先用斑点检测器生成关键点，再用 Lucas--Kanade 跟踪器 \textsuperscript{[707]} 在连续帧之间跟踪这些点。根据跟踪器计算的行驶距离，采用恒速度模型补偿扫描器一次旋转期间的运动，并将变换后的关键点与关键帧位姿一起存储到因子图中。和前述方法一样，关键帧根据行驶距离选取。位姿跟踪使用稀疏关键点，而回环检测则从雷达图中提取更稠密的点云。点云提取方式与 TBV-SLAM \textsuperscript{[9]} 类似；但 RadarSLAM 不选择每个方位角上最强的 $k$ 个点，而是选择强度高于均值一个标准差的所有点。随后从关键帧点云创建 M2DP 描述子 \textsuperscript{[445]}，并通过匹配描述子相似的帧检测回环。为避免匹配缺乏描述性的点云，RadarSLAM 不从过于细长的关键帧中选择回环；也就是说，当 PCA 计算的两个特征值差异明显时，不选择该关键帧，因为这类点云很可能来自高速公路路段等不具唯一性的地点。系统不再执行其他回环验证。

Schuster 等人 \textsuperscript{[983]} 给出了一个基于二维 SoC 雷达的 SLAM 流水线示例。与上述方法不同，该系统维护的图不仅包含传感器位姿，还包含单独的雷达特征节点；TBV-SLAM 和 RadarSLAM 只优化位姿图（尽管会将雷达点关联到位姿节点，以便地点识别和地图渲染）。由于二维 SoC 雷达通常比 360° 旋转雷达或 3+1D SoC 雷达提供少得多的检测结果，因此在可优化图中维护全部检测结果更为可行。图中添加边以表示所有同时观测到的特征之间的观测关系。系统使用二值环状统计描述子（BASD \textsuperscript{[908]}）提取路标。BASD 是紧凑的二值描述子，可以直接比较局部区域内所有特征点的描述子，从而将近期特征与地图中已有特征关联起来。通过 RANSAC 离群点剔除阶段的特征会与带有轮速编码器里程计信息的位姿节点一起作为顶点加入图中。假设开环跟踪的漂移适中，系统不包含显式地点识别步骤；但回环闭合后，可以用 BASD 描述子匹配点特征，随后在后端优化 SLAM 图。

近期两个基于 3+1D SoC 雷达的六自由度雷达 SLAM 方法是 4DRadarSLAM \textsuperscript{[1271]} 和 4D iRIOM \textsuperscript{[1310]}。二者都以三维雷达点云为输入，点检测（对数据立方体进行滤波）由传感器自身在板上完成，因此 SLAM 流水线不会显式包含 CFAR 或其他峰值检测步骤。不过，输入点云仍可能包含大量噪声点。两种方法都利用 3+1D 点云提供的多普勒径向速度信息过滤运动物体的点。车辆自身速度可通过对测得的多普勒点速度执行线性最小二乘估计；对于速度模型不一致的离群点则将其移除。iRIOM 进一步对点云去噪：仅保留满足以下条件的内点，即在固定半径内具有足够多的邻点，且这些邻点的空间分布较为紧凑（根据其空间分布的协方差矩阵判断）。多普勒自身速度估计被用作扫描到子地图配准的先验。4DRadarSLAM 使用一种称为 APDGICP 的 GICP 变体 \textsuperscript{[985]}，通过协方差矩阵为点加权；由于方位角和俯仰角精度有限，距离传感器越远的点具有越高的不确定性。子地图（关键帧）被插入图中。系统通过匹配 Scan Context 检测回环；与上述二维方法不同，它可以使用包含点高度的原始 Scan Context。4D iRIOM 使用原始 Scan Context，而 4DRadarSLAM 使用 Intensity Scan Context \textsuperscript{[1151]}，以避免高度测量噪声导致描述子缺少信息。4DRadarSLAM 还包含验证步骤：若两帧之间累积的里程计距离超过某个阈值，则拒绝 Scan Context 返回的候选匹配。

\subsection{雷达 SLAM 中的多模态}

到目前为止，我们主要讨论只使用雷达数据（有时还使用 IMU 或轮速里程计等本体感知传感器）的雷达 SLAM 方法。本节讨论将毫米波雷达与其他外感知传感器（如 LiDAR、相机）或外部数据（如卫星图像或先验地图）组合的系统。

一些雷达建图工作采用雷达--LiDAR 融合，以根据当前可见性条件生成尽可能好的点集（在清晰条件下更信任 LiDAR，在低能见度条件下更信任雷达）。Fritsche 等人 \textsuperscript{[342, 343]} 根据估计距离融合两种传感器，从而决定在特定情况下信任哪一种。雷达和 LiDAR 也被组合用于改善地点识别：通过将雷达配准到 LiDAR 地图，克服不同传感器模态之间的差异 \textsuperscript{[1247, 1248]}。

Doer 和 Trommer \textsuperscript{[273]} 扩展了 ROVIO \textsuperscript{[90]}（一种基于滤波器的视觉--惯性里程计方法），使用第 9.2.1 节所述的多普勒里程计方法 \textsuperscript{[272]} 集成雷达自身速度估计。类似地，也可以将热成像相机数据与雷达融合，构建多模态雷达--热成像估计流水线 \textsuperscript{[273]}。Zhang 等人 \textsuperscript{[1272]} 将热成像相机数据与 3+1D 雷达点云结合，在低能见度环境中实现鲁棒的帧间里程计。基于 Transformer 的特征匹配器在连续热成像帧中检测对应点，雷达点云则用于改善深度估计。

关于利用多模态数据进行基于雷达的导航，还应提及使用俯视图像和道路地图的方法；不过，文献中的大多数工作专门将这类数据用于定位，而非完整 SLAM。Hong 等人 \textsuperscript{[475]} 展示了如何利用二维扫描雷达数据在 OpenStreetMap 等先验公共地图中定位。其系统使用 RadarSLAM \textsuperscript{[474]} 运行里程计，并将传感器估计位姿表示为高斯混合模型。从 OpenStreetMap 提取代表建筑物墙体的线段作为先验地图特征，再将从雷达图提取的定向点与先验特征匹配。然而，先验地图信息可能不确定、不完整或过时，因此点到特征的数据关联具有挑战性。系统从当前位姿估计的高斯混合模型中采样多个位姿，并在每个位姿下将当前雷达扫描的定向点与先验地图的线特征匹配，从而将扫描定位到先验地图。另一种利用先验地图数据（此处为卫星图像）定位二维雷达图的方法是 RSL-Net \textsuperscript{[1069]}，它由一组深度神经网络组成。第一个网络从输入的俯视照片生成合成图像，展示该地点可能产生的雷达图；另一个网络利用前一网络生成的合成雷达图，估计真实雷达图与俯视图像之间的相对旋转；最后一个网络估计雷达图与俯视图像之间的相对平移偏移。

一些系统使用所谓的超宽带（ultra-wideband，UWB）无线电感知。虽然 UWB 同样使用无线电频谱内的电磁波，有时也被称为 UWB 雷达，但 UWB 的测距能力与毫米波雷达大不相同。在雷达 SLAM 框架中，UWB 雷达主要用于检测不同地点传感器读数之间的相似性。发送宽带、宽波束信号后得到的频率响应可以提供当前位置的特征。度量信息则由轮速编码器或 IMU 数据提供。Schouten 和 Steckel \textsuperscript{[979]} 以及 Takeuchi 等人 \textsuperscript{[1065]} 使用 UWB 波形特征数据库检测重访地点：对于每个地点（节点），存储雷达回波的特征及估计位姿，然后利用里程计和回环约束创建并优化图。两种方法都依靠里程计获得节点之间的距离（即图中的边）来实现度量 SLAM，特征则由脉冲回波 UWB 传感器生成。相反，Premachandra 等人 \textsuperscript{[886]} 使用 UWB 雷达检测点特征，将其用于基于路标的 SLAM 框架。他们在机器人两侧分别安装多个雷达模块，并对两侧传感器信号中匹配峰值进行三边测量，以检测路标。此外，一些基于 UWB 的定位方法采用锚点--标签传感器配置：锚点固定在已知位置，电池供电的 UWB 标签安装在待定位机器人或资产上。但由于这种方法需要预先安装基础设施，它与 SLAM 并不直接相关。

\section{雷达数据集}

本节简要总结雷达 SLAM 文献中具有代表性的数据集，数据集也列于表 9.1。

\textbf{旋转雷达。} 表 9.1 中使用旋转雷达的数据集均采用 Navtech 二维雷达传感器。Oxford Radar RobotCar 数据集 \textsuperscript{[54]} 和 MulRan \textsuperscript{[567]} 是最早用于里程计和 SLAM 的两个大规模雷达数据集，均已被文献广泛使用。Oxford 数据集覆盖城市驾驶路线的多次行驶，总里程达 280 km，包含多种天气条件。MulRan 数据集覆盖更多样的环境，包括密集城市和较偏远的道路行驶，并且不同会话之间的时间跨度更长，但总行驶里程较少。MulRan 主要用于促进地点识别研究，但也广泛用于里程计和 SLAM 方法基准测试。Boreas 数据集 \textsuperscript{[128]} 包含一年内反复行驶同一路线的数据（总计 385 km），尤其包括雪、雨等恶劣天气。除 SLAM 相关基准外，该数据集还包含目标检测基准（汽车、行人和骑行者）。与表 9.1 中其他数据集不同，Oxford Offroad Radar Dataset \textsuperscript{[353]} 关注非城市驾驶，覆盖约 154 km 的非铺装道路和无人居住区域的山间小道。

\textbf{SoC 雷达。} 目前也有若干 SoC 雷达数据集，同时提供 2+1D 和 3+1D 数据。其中一些面向自动驾驶的数据集主要关注目标检测，但也被用于开发和测试 SLAM 方法。nuScenes \textsuperscript{[137]} 将采集车辆上安装的五个 Continental ARS408-21 雷达与一个 LiDAR 和六个相机结合，重点覆盖四座城市的城市驾驶，并且为 23 个目标类别提供了目标检测标注。RadarScenes \textsuperscript{[982]} 是一个包含四个 77 GHz 汽车 2+1D 雷达（未命名）和一个相机的数据集，侧重语义感知，含有 11 类目标的标注，但缺乏准确真值以及 IMU 和 LiDAR 数据。ColoRadar \textsuperscript{[607]} 是一个雷达 SLAM 数据集，包含 3+1D Texas Instruments MMWCAS-RF-EVM 开发板和 2+1D Texas Instruments 模块的数据，以及 IMU 和三维 LiDAR 数据。值得注意的是，该数据集除雷达传感器的原始模数转换器（analog-to-digital converter，ADC）数值、三维“热图”（数据立方体）和单个目标点外，还覆盖室内、室外及地下矿井环境，并包含用于位姿估计的六自由度（6-DoF）真值跟踪。NTU4DRadLM \textsuperscript{[1273]} 和 MSC-RAD4R \textsuperscript{[212]} 均包含 Oculii Eagle 传感器的高分辨率 3+1D 雷达数据。NTU4DRadLM 覆盖结构化（大学校园）和非结构化（公园）环境；MSC-RAD4R 覆盖城市和乡村道路驾驶。Snail-Radar 数据集 \textsuperscript{[482]} 具有两个高分辨率 3+1D 雷达，即 Oculii Eagle 和 Continental ARS548，并包含手持采集和城市道路驾驶数据。

\small
\setlength{\tabcolsep}{2pt}
\begin{center}
\begin{tabular}{p{0.158\linewidth}p{0.158\linewidth}p{0.158\linewidth}p{0.158\linewidth}p{0.158\linewidth}p{0.158\linewidth}}
\toprule
数据集 & LiDAR & 相机 & 真值 & 环境 & 恶劣天气 \\
\midrule
Oxford Radar RobotCar \textsuperscript{[54]} & 有 & 双目/单目 & GPS/IMU + VO & 密集城市 & 雨 \\
Boreas \textsuperscript{[128]} & 有 & 单目 & GPS/IMU + RTK & 稀疏城市 & 雨、雪、雾 \\
MulRan \textsuperscript{[567]} & 有 & 无 & SLAM & 混合城市 & -- \\
RADIATE \textsuperscript{[1001]} & 有 & 双目 & GPS/IMU & 混合城市 & 雨、雪 \\
OORD \textsuperscript{[353]} & 有 & 单目 & GPS & 城市与越野 & 雪、夜间 \\
nuScenes \textsuperscript{[137]} & 有 & 双目 & GPS/IMU & 混合城市与自然环境 & 雨 \\
RadarScenes \textsuperscript{[982]} & 无 & 单目 & 无 & 混合城市道路 & 雨、雾 \\
ColoRadar \textsuperscript{[607]} & 有 & 无 & SLAM & 多种环境 & -- \\
NTU4DRadLM \textsuperscript{[1273]} & 有 & 单目 & SLAM & 混合城市 & -- \\
MSC-RAD4R \textsuperscript{[212]} & 有 & 双目 & GPS + RTK & 城市与乡村道路 & 雪、夜间 \\
Snail \textsuperscript{[482]} & 有 & 双目 & TLS & 道路与隧道 & 雨、夜间 \\
K-Radar \textsuperscript{[836]} & 有 & 双目 & GPS/IMU + RTK & 道路 & 雨、雪、夜间、雾 \\
TruckScene \textsuperscript{[324]} & 有 & 双目 & GPS/IMU + RTK & 道路 & 雨、雪、夜间、雾 \\
HeRCULES \textsuperscript{[569]} & 有 & 双目 & GPS/IMU + RTK & 混合城市与自然环境 & 雨、雪、夜间 \\
\bottomrule
\end{tabular}
\end{center}

表 9.1 公开雷达相关数据集概览。“真值”列中，VO 表示视觉里程计，TLS 表示测量级地面激光扫描，RTK 表示带实时动态校正的 GPS。

\section{延伸阅读与近期趋势}

目前能够处理三维数据的雷达 SLAM 流水线仍然很少，但随着高分辨率 SoC 传感器的发展，可以预期出现更多基于雷达的全三维里程计和地点识别研究。例如，对无人机而言，这一点尤其重要，因为雷达目前在无人机上的应用仍较少。要实现三维大视场雷达单元，需要具有创造性的设计。同时，SLAM 对多个小视场雷达组成的传感器阵列的处理能力也很可能取得显著进展；这自然会伴随标定及其他问题。

研究重点还可能转向更充分地利用包含频谱数据的原始雷达图和数据立方体。传统雷达 SLAM 系统通常通过生成点云来解释环境。然而，未来系统预计会更充分地利用雷达数据中丰富的频谱信息，提供更细致、更有层次的地图，并改进目标检测和分类。其中一个挑战是说服制造商开放其产品原始输出的访问权限，供研究使用。

雷达语义分割也可能在地点识别中发挥更重要的作用。借助分割技术，SLAM 系统能够更有效地区分环境中的不同目标类型，从而实现更智能的导航和决策。这还有助于减少雷达回波中的歧义，使建图更加可靠。

最后，多模态数据方法将受到更多关注。这类方法把雷达数据与其他信息源（包括其他传感器和地理先验，例如 OpenStreetMap）结合起来。例如，将雷达观测与这些先验结合，雷达 SLAM 系统可能能够在大规模室外环境中更准确地定位，进一步提高自主系统的鲁棒性和可靠性。

\textsuperscript{1} 尽管 TX 和 RX 天线有时可以是同一个单元，但多数 SoC 传感器选择将二者分开，从而能够集成多个 TX 和 RX 天线。
\textsuperscript{2} 这些滤波技术可以视为与第 9.2.3 节中的关键点提取类似，但不提取描述子。
\textsuperscript{3} 更全面的概述可参见文献 \textsuperscript{[121]}。
